\chapter{总结与展望}

在实验室能量为72, 68, 64, 60, 56, 52~MeV共六个能点下，首次测量了$^{18}\mathrm{O}+^{152}\mathrm{Sm}$体系的准弹性散射角分布以及弹性散射角分布。其中，准弹性散射角分布至少包含了$^{152}\mathrm{Sm}(^{18}\mathrm{O},^{18}\mathrm{O})^{152}\mathrm{Sm}_{0^+}$、$^{152}\mathrm{Sm}(^{18}\mathrm{O},^{18}\mathrm{O})^{152}\mathrm{Sm}_{2^+}$、$^{152}\mathrm{Sm}(^{18}\mathrm{O},^{18}\mathrm{O})^{152}\mathrm{Sm}_{4^+}$三个反应道，而弹性散射角分布在各个角度的计数是通过多高斯函数在分辨不足的谱中拟合得到的。弹性散射角分布的结果中观察到了明显的库仑虹抑制，这个结果表明，在$^{18}\mathrm{O}+^{152}\mathrm{Sm}$的反应体系中存在大形变导致的强库仑激发的耦合，导致弹性散射反应道被吸收而异常地快速下降。

基于准弹性散射和弹性散射的实验数据，进行了光学模型参数的拟合，提取出各个能点下光学势的实部和虚部的势深度。这也是首次在大形变体系中得到了光学势的能量相依性。结果表明，体系的实部和虚部之间的行为不符合色散关系，还出现了随着能量降低而虚部抬升的行为，这是首次在破裂反应体系之外观察到了反常阈异常的现象。破裂和库仑激发这两种截然不同的反应机制都能与反常阈异常的现象相联系，并且两个体系光学势的变化趋势也有相似性。这些结果与强吸收会影响色散关系成立的猜测一致。考虑到角分布与光学势的联系非常直接，也可以猜测类似观察到角分布抑制的体系，例如$^{11}\mathrm{Li}+^{208}\mathrm{Pb}$或者是$^{16}\mathrm{O}+^{184}\mathrm{W}$的体系也会在光学势的能量相依性上表现出反常阈异常的现象。

值得指出的是，本次实验的能点较少，并且数据的精度还未达到最理想的情况，需要进行更系统、更高精度的实验，才能更完整地描述体系的光学势的行为。这些工作可以对大形变体系的光学势研究提供参考。
